%% ---- sezione B: fatti-stilizzati ----
\section{La matrice di correlazione e i cinque fatti stilizzati}

La matrice di correlazione empirica nasce dai log-rendimenti giornalieri $r_{i,t} = \ln(P_{i,t}/P_{i,t-1})$, preferiti ai prezzi perché stazionari, sui quali si stima la correlazione di Pearson fra coppie di titoli (\S5.2). Essa è simmetrica, a diagonale unitaria, con coefficienti in $[-1, 1]$ e semidefinita positiva (PSD); con finestre di stima lunghe il doppio del numero di titoli (nella tesi $T_w = 2N$) è anche definita positiva, condizione più forte richiesta dalla decomposizione di Cholesky (\S2.3, \S5.2). La validità formale, però, non significa realismo: anche la matrice di correlazione di serie di puro rumore soddisfa questi vincoli. Per giudicare il realismo il capitolo 3 ricorre a cinque fatti stilizzati --- regolarità empiriche persistenti nel senso di Cont (2001) --- già impiegati come protocollo di validazione da Marti (2020) per CorrGAN e da Caprioli et al.\ (2024) per un VAE. Le proprietà sono illustrate sulla matrice di esempio n.~100 ($N = 362$, finestra di 724 giorni).

\begin{itemize}[leftmargin=*, itemsep=2pt]
\item \textbf{Distribuzione pairwise spostata verso valori positivi.} Sotto indipendenza i coefficienti sarebbero centrati sullo zero; nei mercati reali la media è positiva, tipicamente in $[0.2, 0.5]$, per l'effetto mercato. Matrice n.~100: media 0.2334, mediana 0.2257.
\item \textbf{Spettro oltre la banda di Marchenko-Pastur.} Per dati i.i.d.\ nel limite $N, T \to \infty$ con $q = N/T$ fissato, gli autovalori restano nel bulk $[\lambda_-, \lambda_+]$, $\lambda_\pm = (1 \pm \sqrt{q})^2$, e per $q = 0.5$ vale $\lambda_+ = 2.914$. La matrice n.~100 ne ha 9 sopra la soglia: $\lambda_1 = 88.753$, pari al 24.5\% della varianza totale (market mode), e altri otto letti come modi settoriali.
\item \textbf{Proprietà di Perron-Frobenius empirica.} $\lambda_1$ ha molteplicità uno e le componenti dell'autovettore $\mathbf{v}_1$ hanno tutte lo stesso segno. Il teorema fornisce solo una condizione sufficiente (elementi non negativi) che le matrici reali non soddisfano: l'uniformità di segno è una regolarità da verificare caso per caso. Essa è tanto più notevole perché Hüttner e Mai (2019) dimostrano che la possiede solo una frazione $2^{1-d}$ delle matrici di correlazione di dimensione $d$. Una minoranza delle finestre storiche la viola per un singolo titolo debolmente anticorrelato con il resto del paniere (filtro di \S5.3).
\item \textbf{Struttura gerarchica.} Con la distanza di Mantegna $d_{ij} = \sqrt{2(1 - C_{ij})}$ e un clustering ad average linkage, il dendrogramma mostra blocchi corrispondenti ai settori GICS: sulla matrice n.~100 le fusioni intra-settore avvengono tipicamente sotto 1.0, quelle inter-settore sopra 1.15.
\item \textbf{MST scale-free.} Il Minimum Spanning Tree sulle stesse distanze (Mantegna, 1999) ha una distribuzione di grado a legge di potenza $P(k) \propto k^{-\gamma}$ con $2 < \gamma < 3$ (pochi hub, ampia periferia): sulla matrice n.~100, $\gamma = 2.207$ e grado massimo 14.
\end{itemize}

\paragraph{Il benchmark di rumore (\S3.7).} Per verificare che le proprietà non siano artefatti della dimensionalità campionaria, lo stesso protocollo è applicato a una matrice di controllo ottenuta da $N = 362$ serie di rumore gaussiano i.i.d.\ di lunghezza $T = 724$, le stesse dimensioni della matrice reale. Il controllo non supera nessuna delle cinque verifiche: media pairwise 0.0000 con deviazione standard 0.0373, in accordo con il valore atteso $1/\sqrt{T}$; $\lambda_1 = 2.907$ e nessun autovalore oltre $\lambda_+$; 50.8\% di componenti negative in $\mathbf{v}_1$; nessun blocco nel dendrogramma, con fusioni schiacciate a circa $\sqrt{2}$ e colori settoriali interlacciati; grado massimo dell'MST 6 contro 14 e distribuzione di grado concava, senza legge di potenza.

I cinque fatti sono quindi la firma geometrica che separa l'informazione di mercato dal rumore: è ciò che il modello generativo dovrà riprodurre e che, nel capitolo 8, diventa il protocollo di validazione dei 1000 scenari sintetici. Le ricostruzioni del capitolo 7 sono invece valutate con metriche di errore e di similarità fra MST (\S6.5).
